% Entropy gap suppression and Hankel collapse bounds.

\begin{theorem}[熵缺口对非零谱指纹的指数抑制]\label{thm:xi-entropy-gap-exponential-suppression-nonzero-fingerprint}
沿用 \eqref{eq:xi-comoving-scan-fourier} 的有限缺陷口径，并取采样步长 \(\Delta=1\)，令
\(
u_n:=\widehat H(n)
\)
为离散谱指纹序列。
记整数缺陷指数 \(\kappa_{\mathrm{ind}}:=\sum_{j=1}^{\kappa}m_j\) 与缺陷熵 \(S_{\mathrm{def}}\)（\eqref{eq:xi-defect-entropy-closed-form}），并定义熵缺口
\[
\kappa_{\mathrm{ind}}-S_{\mathrm{def}}
=\sum_{j=1}^{\kappa}\frac{m_j}{1+\delta_j}
\qquad\text{（见 \eqref{eq:xi-defect-entropy-index-decomposition}）}.
\]
则对每个整数 \(n\ge 1\) 成立显式上界
\[
|u_n|
\ \le\
4\pi\,e^{-n}\,(\kappa_{\mathrm{ind}}-S_{\mathrm{def}}).
\]
\end{theorem}
\begin{proof}
由 \eqref{eq:xi-comoving-scan-fourier}（取 \(\Delta=1\)）可写
\[
u_n
=4\pi\sum_{j=1}^{\kappa}m_j\,\frac{\delta_j}{1+\delta_j}\,e^{-(1+\delta_j)n}\,e^{-\mathrm{i}\gamma_j n}.
\]
取模并分离指数因子得
\[
|u_n|
\le
4\pi\,e^{-n}\sum_{j=1}^{\kappa}m_j\,\frac{\delta_j}{1+\delta_j}\,(e^{-\delta_j})^{n}.
\]
对任意 \(\delta>0\) 有 \(e^\delta\ge 1+\delta\)，从而 \(e^{-\delta}\le (1+\delta)^{-1}\)；
于是对 \(n\ge 1\),
\(
(e^{-\delta_j})^{n}\le (1+\delta_j)^{-n}\le (1+\delta_j)^{-1}
\)。
代回并用 \(\frac{\delta_j}{1+\delta_j}\le 1\) 得
\[
|u_n|
\le
4\pi\,e^{-n}\sum_{j=1}^{\kappa}m_j\,\frac{1}{1+\delta_j}
=
4\pi\,e^{-n}\,(\kappa_{\mathrm{ind}}-S_{\mathrm{def}}),
\]
其中最后一步使用 \eqref{eq:xi-defect-entropy-index-decomposition}。证毕。
\end{proof}
\leanverified{paper\_xi\_entropy\_gap\_exponential\_suppression\_nonzero\_fingerprint}

\begin{proposition}[非零谐波的熵乘积锐化界]\label{prop:xi-nonzero-harmonic-entropy-gap-product-sharp}
在定理 \ref{thm:xi-entropy-gap-exponential-suppression-nonzero-fingerprint} 的设定下，令
\[
w_j:=\frac{\delta_j}{1+\delta_j}\in(0,1),
\qquad
S_{\mathrm{def}}=\sum_{j=1}^{\kappa}m_j w_j,
\qquad
\kappa_{\mathrm{ind}}-S_{\mathrm{def}}=\sum_{j=1}^{\kappa}m_j(1-w_j).
\]
则对每个整数 \(n\ge 1\) 成立统一上界
\[
|u_n|
\ \le\
4\pi\,e^{-n}\,\frac{S_{\mathrm{def}}(\kappa_{\mathrm{ind}}-S_{\mathrm{def}})}{\kappa_{\mathrm{ind}}}.
\]
\end{proposition}
\begin{proof}
由 \eqref{eq:xi-comoving-scan-fourier}（取 \(\Delta=1\)）并引入 \(w_j=\delta_j/(1+\delta_j)\) 得
\[
u_n
=4\pi\,e^{-n}\sum_{j=1}^{\kappa}m_j\,w_j\,e^{-\delta_j n}\,e^{-\mathrm{i}\gamma_j n}.
\]
由 \(e^\delta\ge 1+\delta\) 推出 \(e^{-\delta}\le (1+\delta)^{-1}=1-w\)，从而对 \(n\ge 1\) 有
\(
e^{-\delta_j n}\le (1-w_j)^n\le (1-w_j)
\)。
因此
\[
|u_n|
\le
4\pi\,e^{-n}\sum_{j=1}^{\kappa}m_j\,w_j(1-w_j)
=
4\pi\,e^{-n}\Bigl(S_{\mathrm{def}}-\sum_{j=1}^{\kappa}m_j w_j^2\Bigr).
\]
将 \((w_j)\) 按重数展开为长度 \(\kappa_{\mathrm{ind}}\) 的序列并应用 Cauchy--Schwarz 不等式，
\[
\sum_{j=1}^{\kappa}m_j w_j^2
\ \ge\
\frac{\bigl(\sum_{j=1}^{\kappa}m_j w_j\bigr)^2}{\sum_{j=1}^{\kappa}m_j}
=\frac{S_{\mathrm{def}}^2}{\kappa_{\mathrm{ind}}}.
\]
代回即得结论。证毕。
\end{proof}
\leanverified{paper\_xi\_nonzero\_harmonic\_entropy\_gap\_product\_sharp}

\begin{corollary}[两点采样反演的熵缺口下界]\label{cor:xi-entropy-gap-lower-bound-from-two-samples}
在命题 \ref{prop:xi-nonzero-harmonic-entropy-gap-product-sharp} 的设定下，若 \(u_0=4\pi S_{\mathrm{def}}\)，则对每个整数 \(n\ge 1\) 有
\[
\kappa_{\mathrm{ind}}-S_{\mathrm{def}}
\ \ge\
\kappa_{\mathrm{ind}}\,e^{n}\,\frac{|u_n|}{|u_0|}.
\]
\end{corollary}
\begin{proof}
由命题 \ref{prop:xi-nonzero-harmonic-entropy-gap-product-sharp} 得
\[
|u_n|
\le
4\pi\,e^{-n}\,\frac{S_{\mathrm{def}}(\kappa_{\mathrm{ind}}-S_{\mathrm{def}})}{\kappa_{\mathrm{ind}}}.
\]
结合 \(u_0=4\pi S_{\mathrm{def}}\) 并整理即得。证毕。
\end{proof}
\leanverified{paper\_xi\_entropy\_gap\_lower\_bound\_from\_two\_samples}

\begin{proposition}[第一谐波的方差惩罚律]\label{prop:xi-first-harmonic-variance-penalty}
\leanverified{paper\_xi\_first\_harmonic\_variance\_penalty}
在命题 \ref{prop:xi-nonzero-harmonic-entropy-gap-product-sharp} 的设定下，将 \((w_j)\) 按重数展开为长度 \(\kappa_{\mathrm{ind}}\) 的序列 \((w_\ell)_{\ell=1}^{\kappa_{\mathrm{ind}}}\subset(0,1)\) 并记均值
\[
\bar w:=\frac{1}{\kappa_{\mathrm{ind}}}\sum_{\ell=1}^{\kappa_{\mathrm{ind}}}w_\ell=\frac{S_{\mathrm{def}}}{\kappa_{\mathrm{ind}}},
\qquad
\Var(w):=\frac{1}{\kappa_{\mathrm{ind}}}\sum_{\ell=1}^{\kappa_{\mathrm{ind}}}(w_\ell-\bar w)^2.
\]
则第一谐波满足显式上界
\[
\frac{e}{4\pi}\,|u_1|
\ \le\
\frac{S_{\mathrm{def}}(\kappa_{\mathrm{ind}}-S_{\mathrm{def}})}{\kappa_{\mathrm{ind}}}
\ -\
\kappa_{\mathrm{ind}}\,\Var(w).
\]
\end{proposition}
\begin{proof}
由 \eqref{eq:xi-comoving-scan-fourier}（取 \(\Delta=1\) 且 \(n=1\)）并沿用 \(e^{-\delta}\le (1+\delta)^{-1}=1-w\) 得
\[
|u_1|
\le
4\pi\,e^{-1}\sum_{\ell=1}^{\kappa_{\mathrm{ind}}}w_\ell(1-w_\ell)
=
4\pi\,e^{-1}\Bigl(\sum_{\ell=1}^{\kappa_{\mathrm{ind}}}w_\ell-\sum_{\ell=1}^{\kappa_{\mathrm{ind}}}w_\ell^2\Bigr).
\]
注意
\[
\sum_{\ell=1}^{\kappa_{\mathrm{ind}}}w_\ell^2
=
\kappa_{\mathrm{ind}}\Var(w)+\kappa_{\mathrm{ind}}\bar w^2
=
\kappa_{\mathrm{ind}}\Var(w)+\frac{S_{\mathrm{def}}^2}{\kappa_{\mathrm{ind}}},
\]
代回并整理即得结论。证毕。
\end{proof}

\begin{proposition}[熵缺口驱动的 Hankel 主子式塌缩上界]\label{prop:xi-entropy-gap-hankel-determinant-collapse-upper-bound}
\leanverified{paper\_xi\_entropy\_gap\_hankel\_determinant\_collapse\_upper\_bound}
令 \(\mathsf H_{\kappa-1}:=(u_{p+q})_{0\le p,q\le \kappa-1}\) 为 Toeplitz--扫描 Hankel 框架中的 Hankel 指纹块，并定义
\[
B:=\sup_{n\ge 1}|u_n|.
\]
则行列式满足组合学型上界
\[
|\det(\mathsf H_{\kappa-1})|
\ \le\
(\kappa-1)!\,|u_0|\,B^{\kappa-1}+\bigl(\kappa!-(\kappa-1)!\bigr)\,B^{\kappa}
\ \le\
\kappa!\bigl(|u_0|\,B^{\kappa-1}+B^{\kappa}\bigr).
\]
特别地，在定理 \ref{thm:xi-entropy-gap-exponential-suppression-nonzero-fingerprint} 的设定下，若 \(S_{\mathrm{def}}\) 固定而 \(\kappa_{\mathrm{ind}}-S_{\mathrm{def}}\to 0\)，则存在仅依赖 \(\kappa\) 的常数使
\[
|\det(\mathsf H_{\kappa-1})|\ \ll_{\kappa}\ (\kappa_{\mathrm{ind}}-S_{\mathrm{def}})^{\kappa-1}.
\]
\end{proposition}
\begin{proof}
以排列展开式写
\[
\det(\mathsf H_{\kappa-1})
=\sum_{\pi\in S_{\kappa}}\mathrm{sgn}(\pi)\prod_{p=0}^{\kappa-1}u_{p+\pi(p)}.
\]
注意 \(p+\pi(p)=0\) 仅可能在 \(p=0\) 且 \(\pi(0)=0\) 时发生，因此每一单项式至多含一个因子 \(u_0\)。
若 \(\pi(0)=0\)，其余因子指标皆满足 \(p+\pi(p)\ge 1\)，故
\[
\Bigl|\prod_{p=0}^{\kappa-1}u_{p+\pi(p)}\Bigr|
\le
|u_0|\,B^{\kappa-1}.
\]
此类排列恰有 \((\kappa-1)!\) 个。
若 \(\pi(0)\neq 0\)，则全部指标 \(\ge 1\)，从而
\[
\Bigl|\prod_{p=0}^{\kappa-1}u_{p+\pi(p)}\Bigr|
\le
B^{\kappa},
\]
此类排列数为 \(\kappa!-(\kappa-1)!\)。
求和得第一条不等式，第二条由粗略放大系数得到。
\par
最后一条由定理 \ref{thm:xi-entropy-gap-exponential-suppression-nonzero-fingerprint} 给出 \(B\ll (\kappa_{\mathrm{ind}}-S_{\mathrm{def}})\)，并结合 \(u_0=4\pi S_{\mathrm{def}}\)（命题 \ref{prop:xi-defect-entropy-closed-form-invariant}）代入即得。证毕。
\end{proof}

\leanverified{paper\_xi\_hankel\_spike\_singular\_spectrum\_separation}
\begin{theorem}[Hankel 主块的一尖峰奇异谱分离]\label{thm:xi-hankel-spike-singular-spectrum-separation}
令 \(\mathsf H_{\kappa-1}:=(u_{p+q})_{0\le p,q\le \kappa-1}\)，并记 \(B:=\sup_{n\ge 1}|u_n|\)。
将 \(\mathsf H_{\kappa-1}\) 的奇异值按降序记为 \(\sigma_1\ge\cdots\ge\sigma_\kappa\)。
则
\[
\sigma_1(\mathsf H_{\kappa-1})\ \ge\ |u_0|-\kappa B,
\qquad
\sigma_j(\mathsf H_{\kappa-1})\ \le\ \kappa B\quad(2\le j\le \kappa).
\]
当 \(\kappa B<|u_0|\) 时，第一奇异值与其余奇异值存在显式分离。
\end{theorem}
\begin{proof}
记标准基向量 \(e_0:=(1,0,\dots,0)^{\mathsf T}\in\RR^{\kappa}\)，并分解
\[
\mathsf H_{\kappa-1}=u_0\,e_0e_0^{\mathsf T}+E,
\qquad
E_{pq}:=\begin{cases}
0,&p=q=0,\\
u_{p+q},&p+q\ge 1.
\end{cases}
\]
由定义 \(|E_{pq}|\le B\)，故
\[
\|E\|_{\mathrm{op}}
\le
\|E\|_{\mathrm{F}}
\le
\sqrt{\kappa^2-1}\,B
\le
\kappa B.
\]
秩一矩阵 \(u_0e_0e_0^{\mathsf T}\) 的奇异值为 \((|u_0|,0,\dots,0)\)。
由 Weyl 奇异值扰动不等式得
\[
\sigma_1(\mathsf H_{\kappa-1})\ge |u_0|-\|E\|_{\mathrm{op}}\ge |u_0|-\kappa B,
\qquad
\sigma_j(\mathsf H_{\kappa-1})\le \|E\|_{\mathrm{op}}\le \kappa B\ (j\ge 2).
\]
证毕。
\end{proof}

\leanverified{paper\_xi\_negative\_block\_spectrum\_spike\_and\_shallow\_bound}
\begin{corollary}[负块谱的单尖峰与浅负谱上界]\label{cor:xi-negative-block-spectrum-spike-and-shallow-bound}
令 \(G_{\kappa-1}:=\mathsf H_{\kappa-1}^{\ast}\mathsf H_{\kappa-1}\succeq 0\)。
则 \(G_{\kappa-1}\) 的特征值满足
\[
\lambda_{\max}(G_{\kappa-1})\ \ge\ (|u_0|-\kappa B)^2,
\qquad
\lambda_j(G_{\kappa-1})\ \le\ (\kappa B)^2\quad(2\le j\le \kappa).
\]
因此负块 \(-G_{\kappa-1}\) 的最深负特征值不大于 \(-(|u_0|-\kappa B)^2\)，其余负特征值皆不小于 \(-(\kappa B)^2\)。
\end{corollary}
\begin{proof}
由定义 \(\lambda_j(G_{\kappa-1})=\sigma_j(\mathsf H_{\kappa-1})^2\)，并应用定理 \ref{thm:xi-hankel-spike-singular-spectrum-separation} 即得。证毕。
\end{proof}
\begin{proposition}[浅负谱方向的秩一可消去半径]\label{prop:xi-rankone-flip-radius-shallow-negative}
在推论 \ref{cor:xi-toeplitz-scan-hankel-congruence-decomposition} 的合同正规形中，负方向由 \(-G_{\kappa-1}\) 给出。
记 \(\lambda_{\min}(G_{\kappa-1})\) 为 \(G_{\kappa-1}\) 的最小特征值，并取其单位特征向量 \(v\)。
则存在 Hermitian 秩一扰动
\(
\Delta:=\lambda_{\min}(G_{\kappa-1})\,vv^{\ast}
\)
使得相应负特征值精确平移至 \(0\)，从而负惯性至少下降 \(1\)，并且
\[
\|\Delta\|_{\mathrm{op}}=\lambda_{\min}(G_{\kappa-1})
=\sigma_\kappa(\mathsf H_{\kappa-1})^2
\le
(\kappa B)^2.
\]
若进一步使用 \(B\le 4\pi e^{-1}\,S_{\mathrm{def}}(\kappa_{\mathrm{ind}}-S_{\mathrm{def}})/\kappa_{\mathrm{ind}}\)（由命题 \ref{prop:xi-nonzero-harmonic-entropy-gap-product-sharp} 在 \(n=1\) 的特例推出），则得到由 \((S_{\mathrm{def}},\kappa_{\mathrm{ind}},\kappa)\) 参数化的显式上界。
\end{proposition}
\begin{proof}
对 Hermitian 矩阵的谱分解，沿最浅负方向的单位特征向量 \(v\) 施加秩一扰动 \(\Delta=\lambda_{\min}(G_{\kappa-1})vv^{\ast}\)，可将 \(-\lambda_{\min}(G_{\kappa-1})\) 精确推至 \(0\)，其余特征值保持不变，故负惯性至少下降 \(1\)。
范数恒等式 \(\|\Delta\|_{\mathrm{op}}=\lambda_{\min}(G_{\kappa-1})\) 由秩一算子谱半径立得。
最后由 \(\lambda_{\min}(G_{\kappa-1})=\sigma_\kappa(\mathsf H_{\kappa-1})^2\) 与定理 \ref{thm:xi-hankel-spike-singular-spectrum-separation} 的 \(\sigma_\kappa\le \kappa B\) 得到 \(\|\Delta\|_{\mathrm{op}}\le (\kappa B)^2\)。
若代入命题 \ref{prop:xi-nonzero-harmonic-entropy-gap-product-sharp}（\(n=1\)）即得显式化版本。证毕。
\end{proof}
\leanverified{paper\_xi\_rankone\_flip\_radius\_shallow\_negative}

\leanverified{paper\_xi\_entropy\_gap\_controls\_geometric\_mean\_radius}
\begin{proposition}[熵缺口控制节点模长的加权几何平均]\label{prop:xi-entropy-gap-controls-geometric-mean-radius}
在 \eqref{eq:xi-comoving-scan-fourier} 的有限缺陷口径下取 \(\Delta=1\)，并令
\[
\rho_j:=e^{-(1+\delta_j)}\in(0,e^{-1})\qquad(1\le j\le \kappa).
\]
则按重数加权的几何平均满足显式上界
\[
\Bigl(\prod_{j=1}^{\kappa}\rho_j^{m_j}\Bigr)^{1/\kappa_{\mathrm{ind}}}
\ \le\
e^{-1}\cdot\frac{\kappa_{\mathrm{ind}}-S_{\mathrm{def}}}{\kappa_{\mathrm{ind}}}.
\]
\end{proposition}
\begin{proof}
由 \(e^{\delta}\ge 1+\delta\) 得 \(e^{-\delta}\le (1+\delta)^{-1}\)，从而
\(
\rho_j=e^{-1}e^{-\delta_j}\le e^{-1}(1+\delta_j)^{-1}
\)。
对按重数展开后的 \(\kappa_{\mathrm{ind}}\) 个因子应用 AM--GM 不等式，
\[
\Bigl(\prod_{j=1}^{\kappa}(1+\delta_j)^{-m_j}\Bigr)^{1/\kappa_{\mathrm{ind}}}
\le
\frac{1}{\kappa_{\mathrm{ind}}}\sum_{j=1}^{\kappa}\frac{m_j}{1+\delta_j}
=\frac{\kappa_{\mathrm{ind}}-S_{\mathrm{def}}}{\kappa_{\mathrm{ind}}},
\]
其中最后一步使用 \eqref{eq:xi-defect-entropy-index-decomposition}。
合并即得结论。证毕。
\end{proof}

\endinput
